\documentclass[11pt, a4paper]{article}
\usepackage[a4paper, margin=2cm]{geometry}
\usepackage{fontspec}
\usepackage[portuguese, bidi=basic, provide=*]{babel}
\babelprovide[import, onchar=ids fonts]{portuguese}
\babelprovide[import, onchar=ids fonts]{english}
\babelfont{rm}{Noto Sans}
\usepackage{amsmath}
\usepackage{amssymb}
\usepackage{booktabs}
\usepackage{array}
\usepackage{graphicx}
\usepackage{enumitem}
\usepackage{listings}
\usepackage{xcolor}
\setlist[itemize]{label=-}
\usepackage{hyperref}

\lstset{
    basicstyle=\ttfamily\small,
    backgroundcolor=\color{gray!10},
    frame=single,
    breaklines=true,
    columns=fullflexible
}

\begin{document}

\noindent
\begin{minipage}[c]{0.18\textwidth}
\includegraphics[height=\linewidth]{misc/logo_ufsc.png}
\end{minipage}
\begin{minipage}[c]{0.79\textwidth}
    \textbf{UNIVERSIDADE FEDERAL DE SANTA CATARINA\\
    CAMPUS BLUMENAU}\\[0.15cm]
    Disciplina: Sistemas a Eventos Discretos\\[0.15cm]
    Adrian Gazzani Felsky dos Anjos (21105100)
\end{minipage}

\begin{center}
    \Large{\textbf{2ª Avaliação — Resolução (Grafos)}}
\end{center}

% ============================================================================
% 1ª QUESTÃO - ALGORITMO DE DIJKSTRA
% ============================================================================
\section*{1ª Questão: Algoritmo de Dijkstra}

\subsection*{1. Evolução das Iterações (Origem $s = 0$)}

\begin{center}
\begin{tabular}{cccccccc}
\toprule
\textbf{Iteração} & \textbf{Fechado ($u$)} & \textbf{Vértice 0} & \textbf{Vértice 1} & \textbf{Vértice 2} & \textbf{Vértice 3} & \textbf{Vértice 4} & \textbf{Vértice 5} \\
\midrule
Inicial & -- & $0$ (--) & $\infty$ (--) & $\infty$ (--) & $\infty$ (--) & $\infty$ (--) & $\infty$ (--) \\
1 & $0$ & \textbf{0 (--)} & $\infty$ (--) & $20$ ($0$) & $40$ ($0$) & $30$ ($0$) & $\infty$ (--) \\
2 & $2$ & $0$ (--) & $\infty$ (--) & \textbf{20 (0)} & $40$ ($0$) & $30$ ($0$) & $\infty$ (--) \\
3 & $4$ & $0$ (--) & $112$ ($4$) & $20$ ($0$) & $40$ ($0$) & \textbf{30 (0)} & $50$ ($4$) \\
4 & $3$ & $0$ (--) & $112$ ($4$) & $20$ ($0$) & \textbf{40 (0)} & $30$ ($0$) & $50$ ($4$) \\
5 & $5$ & $0$ (--) & $112$ ($4$) & $20$ ($0$) & $40$ ($0$) & $30$ ($0$) & \textbf{50 (4)} \\
6 & $1$ & $0$ (--) & \textbf{112 (4)} & $20$ ($0$) & $40$ ($0$) & $30$ ($0$) & $50$ ($4$) \\
\bottomrule
\end{tabular}
\end{center}

\subsection*{2. Resultado Final}

\begin{center}
\begin{tabular}{cccc}
\toprule
\textbf{Destino ($v$)} & \textbf{Custo Mínimo ($\text{dist}[v]$)} & \textbf{Antecessor ($\text{pred}[v]$)} & \textbf{Caminho Mínimo} \\
\midrule
$0$ & $0$ & -- & $0$ \\
$1$ & $112$ & $4$ & $0 \to 4 \to 1$ \\
$2$ & $20$ & $0$ & $0 \to 2$ \\
$3$ & $40$ & $0$ & $0 \to 3$ \\
$4$ & $30$ & $0$ & $0 \to 4$ \\
$5$ & $50$ & $4$ & $0 \to 4 \to 5$ \\
\bottomrule
\end{tabular}
\end{center}

\subsection*{3. Validação Experimental (MATLAB)}
\begin{lstlisting}
% Saida do Console MATLAB - Questao 1 (Dijkstra):
Distancias:    [0 112 20 40 30 50]
Antecessores:  [-- 4 0 0 0 4]
\end{lstlisting}
% Descomente abaixo para inserir a figura gerada no MATLAB:
% \begin{center}
%     \includegraphics[width=0.65\textwidth]{fig_q1_dijkstra.png}
% \end{center}

\vspace{0.2cm}

% ============================================================================
% 2ª QUESTÃO - ALGORITMO DE BELLMAN-FORD
% ============================================================================
\section*{2ª Questão: Algoritmo de Bellman-Ford}

\subsection*{1. Evolução das Iterações (Origem $s = A$)}

\begin{center}
\begin{tabular}{cccccccc}
\toprule
\textbf{Iteração ($k$)} & \textbf{A} & \textbf{B} & \textbf{C} & \textbf{D} & \textbf{E} & \textbf{F} & \textbf{G} \\
\midrule
$k = 0$ & $0$ (--) & $\infty$ (--) & $\infty$ (--) & $\infty$ (--) & $\infty$ (--) & $\infty$ (--) & $\infty$ (--) \\
$k = 1$ & $0$ (--) & $2$ ($A$) & $\infty$ (--) & $5$ ($A$) & $\infty$ (--) & $3$ ($A$) & $\infty$ (--) \\
$k = 2$ & $0$ (--) & $-1$ ($F$) & $\infty$ (--) & $5$ ($A$) & $3$ ($B$) & $3$ ($A$) & $\infty$ (--) \\
$k = 3$ & $0$ (--) & $-1$ ($F$) & $0$ ($E$) & $5$ ($A$) & $0$ ($B$) & $3$ ($A$) & $6$ ($E$) \\
$k = 4$ & $0$ (--) & $-1$ ($F$) & $-3$ ($E$) & $5$ ($A$) & $0$ ($B$) & $3$ ($A$) & $3$ ($E$) \\
$k = 5$ & $0$ (--) & $-1$ ($F$) & $-3$ ($E$) & $2$ ($G$) & $0$ ($B$) & $3$ ($A$) & $1$ ($C$) \\
$k = 6$ & \textbf{0 (--)} & \textbf{-1 (F)} & \textbf{-3 (E)} & \textbf{0 (G)} & \textbf{0 (B)} & \textbf{3 (A)} & \textbf{1 (C)} \\
\bottomrule
\end{tabular}
\end{center}

\subsection*{2. Tabela Final de Custos e Roteamento}

\begin{center}
\begin{tabular}{cccc}
\toprule
\textbf{Node} & \textbf{Cost} & \textbf{Previous} & \textbf{Caminho Mínimo} \\
\midrule
$A$ & $0$ & -- & $A$ \\
$B$ & $-1$ & $F$ & $A \to F \to B$ \\
$C$ & $-3$ & $E$ & $A \to F \to B \to E \to C$ \\
$D$ & $0$ & $G$ & $A \to F \to B \to E \to C \to G \to D$ \\
$E$ & $0$ & $B$ & $A \to F \to B \to E$ \\
$F$ & $3$ & $A$ & $A \to F$ \\
$G$ & $1$ & $C$ & $A \to F \to B \to E \to C \to G$ \\
\bottomrule
\end{tabular}
\end{center}

\subsection*{3. Validação Experimental (MATLAB)}
\begin{lstlisting}
% Saida do Console MATLAB - Questao 2 (Bellman-Ford):
Distancias:              [0 -1 -3 0 0 3 1]
Predecessores:           [-- F E G B A C]
Ciclo negativo presente: 0
\end{lstlisting}
% Descomente abaixo para inserir a figura gerada no MATLAB:
% \begin{center}
%     \includegraphics[width=0.65\textwidth]{fig_q2_bellman_ford.png}
% \end{center}

\vspace{0.2cm}

% ============================================================================
% 3ª QUESTÃO - ALGORITMO DE PRIM
% ============================================================================
\section*{3ª Questão: Algoritmo de Prim}

\subsection*{1. Evolução da Seleção de Vértices (Início em $A$)}

\begin{center}
\begin{tabular}{ccccccccc}
\toprule
\textbf{Passo} & \textbf{Extraído ($u$)} & \textbf{A} & \textbf{B} & \textbf{C} & \textbf{D} & \textbf{E} & \textbf{F} & \textbf{G} \\
\midrule
Início & -- & $0$ (--) & $\infty$ (--) & $\infty$ (--) & $\infty$ (--) & $\infty$ (--) & $\infty$ (--) & $\infty$ (--) \\
1 & $A$ & \textbf{0 (--)} & $7$ ($A$) & $\infty$ (--) & $5$ ($A$) & $\infty$ (--) & $\infty$ (--) & $\infty$ (--) \\
2 & $D$ & -- & $7$ ($A$) & $\infty$ (--) & \textbf{5 (A)} & $15$ ($D$) & $6$ ($D$) & $\infty$ (--) \\
3 & $F$ & -- & $7$ ($A$) & $\infty$ (--) & -- & $8$ ($F$) & \textbf{6 (D)} & $11$ ($F$) \\
4 & $B$ & -- & \textbf{7 (A)} & $8$ ($B$) & -- & $7$ ($B$) & -- & $11$ ($F$) \\
5 & $E$ & -- & -- & $5$ ($E$) & -- & \textbf{7 (B)} & -- & $9$ ($E$) \\
6 & $C$ & -- & -- & \textbf{5 (E)} & -- & -- & -- & $9$ ($E$) \\
7 & $G$ & -- & -- & -- & -- & -- & -- & \textbf{9 (E)} \\
\bottomrule
\end{tabular}
\end{center}

\subsection*{2. Árvore Geradora Mínima (AGM)}

\begin{center}
\begin{tabular}{ccc}
\toprule
\textbf{Ordem} & \textbf{Aresta $(\text{pai}[u], u)$} & \textbf{Peso} \\
\midrule
1ª & $(A, D)$ & $5$ \\
2ª & $(D, F)$ & $6$ \\
3ª & $(A, B)$ & $7$ \\
4ª & $(B, E)$ & $7$ \\
5ª & $(E, C)$ & $5$ \\
6ª & $(E, G)$ & $9$ \\
\midrule
\multicolumn{2}{r}{\textbf{Peso Total da AGM:}} & \textbf{39} \\
\bottomrule
\end{tabular}
\end{center}

\subsection*{3. Validação Experimental (MATLAB)}
\begin{lstlisting}
% Saida do Console MATLAB - Questao 3 (Prim):
Arestas da AGM: (A,D)=5, (D,F)=6, (A,B)=7, (B,E)=7, (E,C)=5, (E,G)=9
Peso total: 39
\end{lstlisting}
% Descomente abaixo para inserir a figura gerada no MATLAB:
% \begin{center}
%     \includegraphics[width=0.65\textwidth]{fig_q3_prim.png}
% \end{center}

\vspace{0.2cm}

% ============================================================================
% 4ª QUESTÃO - ALGORITMO DE FORD-FULKERSON
% ============================================================================
\section*{4ª Questão: Algoritmo de Ford-Fulkerson}

\subsection*{1. Caminhos Aumentantes ($s = 0 \to t = 5$)}

\begin{center}
\begin{tabular}{cccc}
\toprule
\textbf{Iteração} & \textbf{Caminho Aumentante} & \textbf{Gargalo ($\Delta$)} & \textbf{Fluxo Acumulado} \\
\midrule
1 & $0 \to 2 \to 4 \to 5$ & $\min(13, 14, 4) = 4$ & $4$ \\
2 & $0 \to 2 \to 4 \to 3 \to 5$ & $\min(9, 10, 7, 20) = 7$ & $11$ \\
3 & $0 \to 1 \to 3 \to 5$ & $\min(16, 12, 13) = 12$ & $\mathbf{23}$ \\
\bottomrule
\end{tabular}
\end{center}

\subsection*{2. Fluxo Final e Capacidade Residual nas Arestas}

\begin{center}
\begin{tabular}{cccc}
\toprule
\textbf{Aresta $(u \to v)$} & \textbf{Capacidade $c(u,v)$} & \textbf{Fluxo Final $f(u,v)$} & \textbf{Capacidade Residual $c_f(u,v)$} \\
\midrule
$0 \to 1$ & $16$ & $12$ & $4$ \\
$0 \to 2$ & $13$ & $11$ & $2$ \\
$1 \to 2$ & $10$ & $0$ & $10$ \\
$2 \to 1$ & $4$ & $0$ & $4$ \\
$1 \to 3$ & $12$ & $12$ & $0$ \\
$3 \to 2$ & $9$ & $0$ & $9$ \\
$2 \to 4$ & $14$ & $11$ & $3$ \\
$4 \to 3$ & $7$ & $7$ & $0$ \\
$3 \to 5$ & $20$ & $19$ & $1$ \\
$4 \to 5$ & $4$ & $4$ & $0$ \\
\midrule
\multicolumn{2}{r}{\textbf{Fluxo Máximo Total ($0 \to 5$):}} & \multicolumn{2}{c}{\textbf{23}} \\
\bottomrule
\end{tabular}
\end{center}

\subsection*{3. Validação Experimental (MATLAB)}
\begin{lstlisting}
% Saida do Console MATLAB - Questao 4 (Ford-Fulkerson):
Fluxo maximo: 23
Fluxos ativos: (0->1)=12, (0->2)=11, (1->3)=12, (2->4)=11, (3->5)=19, (4->3)=7, (4->5)=4
\end{lstlisting}
% Descomente abaixo para inserir a figura gerada no MATLAB:
% \begin{center}
%     \includegraphics[width=0.65\textwidth]{fig_q4_ford_fulkerson.png}
% \end{center}

\vspace{0.2cm}

% ============================================================================
% 5ª QUESTÃO - ALGORITMO DE KRUSKAL
% ============================================================================
\section*{5ª Questão: Algoritmo de Kruskal}

\subsection*{1. Processamento de Arestas Ordenadas}
Rotulação posicional dos vértices: $A$ (Sup. Esquerdo), $B$ (Sup. Central), $C$ (Sup. Direito), $D$ (Inf. Esquerdo), $E$ (Inf. Central), $F$ (Inf. Direito).

\begin{center}
\begin{tabular}{ccccc}
\toprule
\textbf{Passo} & \textbf{Aresta} & \textbf{Peso} & \textbf{Decisão} & \textbf{Conjuntos Disjuntos (Union-Find)} \\
\midrule
Início & -- & -- & -- & $\{A\}, \{B\}, \{C\}, \{D\}, \{E\}, \{F\}$ \\
1 & $(A, E)$ & $1$ & \textbf{Aceita} & $\{A, E\}, \{B\}, \{C\}, \{D\}, \{F\}$ \\
2 & $(A, B)$ & $2$ & \textbf{Aceita} & $\{A, B, E\}, \{C\}, \{D\}, \{F\}$ \\
3 & $(B, C)$ & $3$ & \textbf{Aceita} & $\{A, B, C, E\}, \{D\}, \{F\}$ \\
4 & $(C, E)$ & $3$ & Rejeitada & $\{A, B, C, E\}, \{D\}, \{F\}$ \\
5 & $(D, E)$ & $4$ & \textbf{Aceita} & $\{A, B, C, D, E\}, \{F\}$ \\
6 & $(E, F)$ & $4$ & \textbf{Aceita} & $\{A, B, C, D, E, F\}$ \\
7 & $(A, D)$ & $5$ & Rejeitada & $\{A, B, C, D, E, F\}$ \\
8 & $(B, E)$ & $7$ & Rejeitada & $\{A, B, C, D, E, F\}$ \\
9 & $(C, F)$ & $9$ & Rejeitada & $\{A, B, C, D, E, F\}$ \\
\midrule
\multicolumn{2}{r}{\textbf{Peso Total da AGM:}} & \textbf{14} & \multicolumn{2}{l}{\textbf{AGM:} $\{(A,E), (A,B), (B,C), (D,E), (E,F)\}$} \\
\bottomrule
\end{tabular}
\end{center}

\subsection*{2. Validação Experimental (MATLAB)}
\begin{lstlisting}
% Saida do Console MATLAB - Questao 5 (Kruskal):
Arestas da AGM: (A,E)=1, (A,B)=2, (B,C)=3, (D,E)=4, (E,F)=4
Peso total: 14
\end{lstlisting}
% Descomente abaixo para inserir a figura gerada no MATLAB:
% \begin{center}
%     \includegraphics[width=0.65\textwidth]{fig_q5_kruskal.png}
% \end{center}

\end{document}\